\chapter{Positivity of line bundles}
\label{chap:positivity}

% archon:covers MR4276287UniformityInMordellLangForCurves/Positivity.lean

This chapter sets up the definitions of nef and big line bundles and states Siu's
bigness criterion, which is the key external input used in the proof of
\cref{prop:aux_ht_ineq}.  The notions of nef and big are not available in Mathlib
and are introduced here as project-specific definitions.  The chapter also records
the nefness of the specific line bundles \(\mathcal{F}\) and \(\mathcal{M}\) on the
projective closure \(\overline{X_N}\) arising from the graph construction of
\cref{lem:graph_construction}.

\section{Intersection theory and height inequality on the total space}

\begin{definition}[Nef line bundle]
  \label{def:isNef}
  \lean{MR4276287UniformityInMordellLangForCurves.Positivity.isNef}
  Let \(X\) be a projective variety over a field \(k\) and let \(L\) be a line
  bundle on \(X\).  We say \(L\) is \emph{nef} (numerically effective) if
  \[
    \deg(L|_C) \ge 0
  \]
  for every irreducible curve \(C \subset X\), where \(\deg(L|_C) = (L \cdot C)\)
  denotes the intersection number of \(L\) with the curve class \([C]\).
\end{definition}

\begin{definition}[Big line bundle]
  \label{def:isBig}
  \lean{MR4276287UniformityInMordellLangForCurves.Positivity.isBig}
  Let \(X\) be a projective variety of dimension \(n\) over a field \(k\) and let
  \(L\) be a line bundle on \(X\).  We say \(L\) is \emph{big} if
  \[
    h^0(X,\, L^{\otimes m}) \;\sim\; c \cdot m^n \quad \text{as } m \to \infty
  \]
  for some constant \(c > 0\).  Equivalently (by Kodaira's lemma), \(L\) is big
  if and only if there exist a positive integer \(m_0\) and an effective divisor
  \(E\) on \(X\) such that \(L^{\otimes m_0} \otimes \mathcal{O}_X(-E)\) is ample.
\end{definition}

\begin{lemma}[Non-negativity of nef intersection numbers]
  \label{lem:nefIntersection}
  \lean{MR4276287UniformityInMordellLangForCurves.Positivity.nefIntersection}
  \uses{def:isNef}
  Let \(X\) be a projective variety of dimension \(n\) over a field \(k\), let
  \(H\) be an ample line bundle on \(X\), and let \(L_1, \ldots, L_d\) be nef
  line bundles on \(X\) with \(d \le n\).  Then the intersection number
  \[
    (L_1 \cdot L_2 \cdots L_d \cdot H^{n-d}) \;\ge\; 0.
  \]
  In particular, taking \(d = n\), if \(L_1, \ldots, L_n\) are nef then
  \((L_1 \cdot L_2 \cdots L_n) \ge 0\).
\end{lemma}

\begin{proof}
  \uses{def:isNef}
  Since each \(L_i\) is nef, for every integer \(r \ge 1\) the bundle
  \(L_i^{\otimes r} \otimes H\) is ample (by the numerical characterization of
  nefness in characteristic zero).  Multi-linearity and continuity of the
  intersection pairing then give
  \[
    (L_1 \cdot \cdots \cdot L_d \cdot H^{n-d})
    = \lim_{r \to \infty} r^{-d}\,
      \bigl((L_1^{\otimes r} \otimes H) \cdots (L_d^{\otimes r} \otimes H)
      \cdot H^{n-d}\bigr) \ge 0,
  \]
  where the right-hand side is non-negative because it is an intersection number
  of ample bundles on a projective variety.
\end{proof}

% SOURCE: Section 4 of \cite{MR4276287-mordell-lang-curves}
%         (read from references/MR4276287-mordell-lang-curves-source/HtIneq.tex)
% SOURCE QUOTE: "Both $\cF$ and $\cM$ are nef line bundles. Thus a criterion of
% bigness by Siu \cite[Theorem~2.2.15]{PosAlgGeom}, states that $\cF^{\otimes q}
% \otimes \cM^{\otimes -p N^2}$ is big if $(\cF^{\cdot d}) > d c_1
% (\cM^{\otimes N^2} \cdot \cF^{\cdot (d-1)})$. Note that $(\cM^{\otimes N^2}
% \cdot \cF^{\cdot (d-1)}) = N^2 (\cM \cdot \cF^{\cdot (d-1)})$ by
% multi-linearity of intersection numbers."
\begin{theorem}[Siu's bigness criterion]
  \label{thm:siuCriterion}
  \lean{MR4276287UniformityInMordellLangForCurves.Positivity.siuCriterion}
  \uses{def:isNef, def:isBig, lem:nefIntersection}
  \textit{Source: \cite[Theorem~2.2.15]{PosAlgGeom}, as cited in Section~4 of
  \cite{MR4276287-mordell-lang-curves}
  (read from references/MR4276287-mordell-lang-curves-source/HtIneq.tex).}

  Let \(X\) be a projective variety of dimension \(n\) over a field \(k\), and
  let \(L\) and \(M\) be nef line bundles on \(X\).  Suppose that
  \[
    (L^n) \;>\; n \cdot (M \cdot L^{n-1}).
  \]
  Then there exists a positive integer \(q\) such that
  \(L^{\otimes q} \otimes M^{-1}\) is big.  In particular, there is an effective
  Cartier divisor in the complete linear system
  \(\bigl|L^{\otimes q} \otimes M^{-1}\bigr|\).
\end{theorem}

% SOURCE QUOTE PROOF: Proof of Theorem 2.2.15 is in Lazarsfeld,
% \textit{Positivity in Algebraic Geometry I} (not available in local
% references/; verbatim text not yet retrieved).
\begin{proof}
  \uses{def:isNef, def:isBig, lem:nefIntersection}
  This is \cite[Theorem~2.2.15]{PosAlgGeom}.  The proof constructs a suitable
  \(\mathbb{Q}\)-linear combination of \(L\) and \(M\) that is big, exploiting the
  positivity of the self-intersection \((L^n)\) relative to the mixed intersection
  \((M \cdot L^{n-1})\) together with the non-negativity of nef intersection
  numbers from \cref{lem:nefIntersection}; we refer to Lazarsfeld for details.
  In the application \cref{prop:aux_ht_ineq}, this criterion is invoked with
  \(L = \mathcal{F}^{\otimes q}\), \(M = \mathcal{M}^{\otimes qc_1N^2}\), and
  intersection bounds from the two intersection-number propositions in
  Section~4 of \cite{MR4276287-mordell-lang-curves}.
\end{proof}

% SOURCE: Section 4 of \cite{MR4276287-mordell-lang-curves}
%         (read from references/MR4276287-mordell-lang-curves-source/HtIneq.tex)
% SOURCE QUOTE: "Both $\cF$ and $\cM$ are nef line bundles."
\begin{lemma}[The graph line bundles are nef]
  \label{lem:nefFromBetti}
  \lean{MR4276287UniformityInMordellLangForCurves.Positivity.nefFromBetti}
  \uses{def:isNef, lem:graph_construction}
  \textit{Source: Section~4 of \cite{MR4276287-mordell-lang-curves}
  (read from references/MR4276287-mordell-lang-curves-source/HtIneq.tex).}

  Let \(\overline{X_N} \subset \PP^n \times \PP^n \times \PP^m\) be the
  projective closure from \cref{lem:graph_construction}.  The line bundles
  \[
    \mathcal{F} = \mathcal{O}(0,1,1)\big|_{\overline{X_N}}
    \quad \text{and} \quad
    \mathcal{M} = \mathcal{O}(0,0,1)\big|_{\overline{X_N}}
  \]
  are nef on \(\overline{X_N}\).
\end{lemma}

% SOURCE QUOTE PROOF: "we rely on the following well-known positivity property
% of the intersection theory of multiprojective space: any effective Weil divisor
% on a multiprojective space is nef. This approach was motivated by
% K\"uhne's~\cite{kuehne:semiabelianbhc} work on semiabelian varieties."
% (read from references/MR4276287-mordell-lang-curves-source/HtIneq.tex)
\begin{proof}
  \uses{def:isNef, lem:graph_construction}
  The bundle \(\mathcal{O}(0,0,1)\) on \(\PP^n \times \PP^n \times \PP^m\) is
  the pullback of \(\mathcal{O}(1)\) from the third factor \(\PP^m\); similarly,
  \(\mathcal{O}(0,1,1)\) is the pullback of \(\mathcal{O}(1)\) from the second and
  third factors combined via the Segre embedding.  Pullbacks of globally generated
  line bundles are globally generated, hence nef.

  Alternatively, as noted in the proof of Proposition~4.2 of
  \cite{MR4276287-mordell-lang-curves}, any effective Weil divisor on a
  multiprojective space is nef: because every irreducible curve \(C\) in
  \(\PP^n \times \PP^n \times \PP^m\) meets an effective hyperplane divisor
  in a non-negative intersection number (it cannot be contained in every
  such hyperplane class), the tautological \(\mathcal{O}(0,0,1)\) and
  \(\mathcal{O}(0,1,1)\) are nef on the ambient space.  Since \(\overline{X_N}\)
  is a closed subvariety, the restrictions
  \(\mathcal{O}(0,0,1)|_{\overline{X_N}}\) and
  \(\mathcal{O}(0,1,1)|_{\overline{X_N}}\) are nef on \(\overline{X_N}\).
\end{proof}
